\IfFileExists{stacks-project.cls}{%
\documentclass{stacks-project}
}{%
\documentclass{amsart}
}

% For dealing with references we use the comment environment
\usepackage{verbatim}
\newenvironment{reference}{\comment}{\endcomment}
%\newenvironment{reference}{}{}
\newenvironment{slogan}{\comment}{\endcomment}
\newenvironment{history}{\comment}{\endcomment}

% For commutative diagrams we use Xy-pic
\usepackage[all]{xy}

% We use 2cell for 2-commutative diagrams.
\xyoption{2cell}
\UseAllTwocells

% We use multicol for the list of chapters between chapters
\usepackage{multicol}

% This is generally recommended for better output
\usepackage{lmodern}
\usepackage[T1]{fontenc}

% For cross-file-references
\usepackage{xr-hyper}

% Package for hypertext links:
\usepackage{hyperref}

% For any local file, say "hello.tex" you want to link to please
% use \externaldocument[hello-]{hello}
\externaldocument[introduction-]{introduction}
\externaldocument[conventions-]{conventions}
\externaldocument[sets-]{sets}
\externaldocument[categories-]{categories}
\externaldocument[topology-]{topology}
\externaldocument[sheaves-]{sheaves}
\externaldocument[sites-]{sites}
\externaldocument[stacks-]{stacks}
\externaldocument[fields-]{fields}
\externaldocument[algebra-]{algebra}
\externaldocument[brauer-]{brauer}
\externaldocument[homology-]{homology}
\externaldocument[derived-]{derived}
\externaldocument[simplicial-]{simplicial}
\externaldocument[more-algebra-]{more-algebra}
\externaldocument[smoothing-]{smoothing}
\externaldocument[modules-]{modules}
\externaldocument[sites-modules-]{sites-modules}
\externaldocument[injectives-]{injectives}
\externaldocument[cohomology-]{cohomology}
\externaldocument[sites-cohomology-]{sites-cohomology}
\externaldocument[dga-]{dga}
\externaldocument[dpa-]{dpa}
\externaldocument[sdga-]{sdga}
\externaldocument[hypercovering-]{hypercovering}
\externaldocument[schemes-]{schemes}
\externaldocument[constructions-]{constructions}
\externaldocument[properties-]{properties}
\externaldocument[morphisms-]{morphisms}
\externaldocument[coherent-]{coherent}
\externaldocument[divisors-]{divisors}
\externaldocument[limits-]{limits}
\externaldocument[varieties-]{varieties}
\externaldocument[topologies-]{topologies}
\externaldocument[descent-]{descent}
\externaldocument[perfect-]{perfect}
\externaldocument[more-morphisms-]{more-morphisms}
\externaldocument[flat-]{flat}
\externaldocument[groupoids-]{groupoids}
\externaldocument[more-groupoids-]{more-groupoids}
\externaldocument[etale-]{etale}
\externaldocument[chow-]{chow}
\externaldocument[intersection-]{intersection}
\externaldocument[pic-]{pic}
\externaldocument[weil-]{weil}
\externaldocument[adequate-]{adequate}
\externaldocument[dualizing-]{dualizing}
\externaldocument[duality-]{duality}
\externaldocument[discriminant-]{discriminant}
\externaldocument[derham-]{derham}
\externaldocument[local-cohomology-]{local-cohomology}
\externaldocument[algebraization-]{algebraization}
\externaldocument[curves-]{curves}
\externaldocument[resolve-]{resolve}
\externaldocument[models-]{models}
\externaldocument[functors-]{functors}
\externaldocument[equiv-]{equiv}
\externaldocument[pione-]{pione}
\externaldocument[etale-cohomology-]{etale-cohomology}
\externaldocument[proetale-]{proetale}
\externaldocument[relative-cycles-]{relative-cycles}
\externaldocument[more-etale-]{more-etale}
\externaldocument[trace-]{trace}
\externaldocument[crystalline-]{crystalline}
\externaldocument[spaces-]{spaces}
\externaldocument[spaces-properties-]{spaces-properties}
\externaldocument[spaces-morphisms-]{spaces-morphisms}
\externaldocument[decent-spaces-]{decent-spaces}
\externaldocument[spaces-cohomology-]{spaces-cohomology}
\externaldocument[spaces-limits-]{spaces-limits}
\externaldocument[spaces-divisors-]{spaces-divisors}
\externaldocument[spaces-over-fields-]{spaces-over-fields}
\externaldocument[spaces-topologies-]{spaces-topologies}
\externaldocument[spaces-descent-]{spaces-descent}
\externaldocument[spaces-perfect-]{spaces-perfect}
\externaldocument[spaces-more-morphisms-]{spaces-more-morphisms}
\externaldocument[spaces-flat-]{spaces-flat}
\externaldocument[spaces-groupoids-]{spaces-groupoids}
\externaldocument[spaces-more-groupoids-]{spaces-more-groupoids}
\externaldocument[bootstrap-]{bootstrap}
\externaldocument[spaces-pushouts-]{spaces-pushouts}
\externaldocument[spaces-chow-]{spaces-chow}
\externaldocument[groupoids-quotients-]{groupoids-quotients}
\externaldocument[spaces-more-cohomology-]{spaces-more-cohomology}
\externaldocument[spaces-simplicial-]{spaces-simplicial}
\externaldocument[spaces-duality-]{spaces-duality}
\externaldocument[formal-spaces-]{formal-spaces}
\externaldocument[restricted-]{restricted}
\externaldocument[spaces-resolve-]{spaces-resolve}
\externaldocument[formal-defos-]{formal-defos}
\externaldocument[defos-]{defos}
\externaldocument[cotangent-]{cotangent}
\externaldocument[examples-defos-]{examples-defos}
\externaldocument[algebraic-]{algebraic}
\externaldocument[examples-stacks-]{examples-stacks}
\externaldocument[stacks-sheaves-]{stacks-sheaves}
\externaldocument[criteria-]{criteria}
\externaldocument[artin-]{artin}
\externaldocument[quot-]{quot}
\externaldocument[stacks-properties-]{stacks-properties}
\externaldocument[stacks-morphisms-]{stacks-morphisms}
\externaldocument[stacks-limits-]{stacks-limits}
\externaldocument[stacks-cohomology-]{stacks-cohomology}
\externaldocument[stacks-perfect-]{stacks-perfect}
\externaldocument[stacks-introduction-]{stacks-introduction}
\externaldocument[stacks-more-morphisms-]{stacks-more-morphisms}
\externaldocument[stacks-geometry-]{stacks-geometry}
\externaldocument[moduli-]{moduli}
\externaldocument[moduli-curves-]{moduli-curves}
\externaldocument[examples-]{examples}
\externaldocument[exercises-]{exercises}
\externaldocument[guide-]{guide}
\externaldocument[desirables-]{desirables}
\externaldocument[coding-]{coding}
\externaldocument[obsolete-]{obsolete}
\externaldocument[fdl-]{fdl}
\externaldocument[index-]{index}

% Theorem environments.
%
\theoremstyle{plain}
\newtheorem{theorem}[subsection]{Theorem}
\newtheorem{proposition}[subsection]{Proposition}
\newtheorem{lemma}[subsection]{Lemma}

\theoremstyle{definition}
\newtheorem{definition}[subsection]{Definition}
\newtheorem{example}[subsection]{Example}
\newtheorem{exercise}[subsection]{Exercise}
\newtheorem{situation}[subsection]{Situation}

\theoremstyle{remark}
\newtheorem{remark}[subsection]{Remark}
\newtheorem{remarks}[subsection]{Remarks}

\numberwithin{equation}{subsection}

% Macros
%
\def\lim{\mathop{\mathrm{lim}}\nolimits}
\def\colim{\mathop{\mathrm{colim}}\nolimits}
\def\Spec{\mathop{\mathrm{Spec}}}
\def\Hom{\mathop{\mathrm{Hom}}\nolimits}
\def\Ext{\mathop{\mathrm{Ext}}\nolimits}
\def\SheafHom{\mathop{\mathcal{H}\!\mathit{om}}\nolimits}
\def\SheafExt{\mathop{\mathcal{E}\!\mathit{xt}}\nolimits}
\def\Sch{\mathit{Sch}}
\def\Mor{\mathop{\mathrm{Mor}}\nolimits}
\def\Ob{\mathop{\mathrm{Ob}}\nolimits}
\def\Sh{\mathop{\mathit{Sh}}\nolimits}
\def\NL{\mathop{N\!L}\nolimits}
\def\CH{\mathop{\mathrm{CH}}\nolimits}
\def\proetale{{pro\text{-}\acute{e}tale}}
\def\etale{{\acute{e}tale}}
\def\QCoh{\mathit{QCoh}}
\def\Ker{\mathop{\mathrm{Ker}}}
\def\Im{\mathop{\mathrm{Im}}}
\def\Coker{\mathop{\mathrm{Coker}}}
\def\Coim{\mathop{\mathrm{Coim}}}

% Boxtimes
%
\DeclareMathSymbol{\boxtimes}{\mathbin}{AMSa}{"02}

%
% Macros for moduli stacks/spaces
%
\def\QCohstack{\mathcal{QC}\!\mathit{oh}}
\def\Cohstack{\mathcal{C}\!\mathit{oh}}
\def\Spacesstack{\mathcal{S}\!\mathit{paces}}
\def\Quotfunctor{\mathrm{Quot}}
\def\Hilbfunctor{\mathrm{Hilb}}
\def\Curvesstack{\mathcal{C}\!\mathit{urves}}
\def\Polarizedstack{\mathcal{P}\!\mathit{olarized}}
\def\Complexesstack{\mathcal{C}\!\mathit{omplexes}}
% \Pic is the operator that assigns to X its picard group, usage \Pic(X)
% \Picardstack_{X/B} denotes the Picard stack of X over B
% \Picardfunctor_{X/B} denotes the Picard functor of X over B
\def\Pic{\mathop{\mathrm{Pic}}\nolimits}
\def\Picardstack{\mathcal{P}\!\mathit{ic}}
\def\Picardfunctor{\mathrm{Pic}}
\def\Deformationcategory{\mathcal{D}\!\mathit{ef}}
\usepackage{cleveref}
\crefname{originalreference}{original reference}{original references}
\Crefname{originalreference}{Original reference}{Original references}
\makeatletter
\newlabel{lemma-describe-sheaves-simplicial-site-site}{{09WF}{source}{}{sourcepointer-1}{}}
\newlabel{lemma-describe-sheaves-simplicial-site-site@cref}{{[originalreference][0][]09WF}{[1][0][]source}{}{}{}}
\makeatother
\crefname{originalreference}{original reference}{original references}
\Crefname{originalreference}{Original reference}{Original references}
\makeatletter
\newlabel{lemma-restriction-to-components-site}{{09WG}{source}{}{sourcepointer-2}{}}
\newlabel{lemma-restriction-to-components-site@cref}{{[originalreference][0][]09WG}{[1][0][]source}{}{}{}}
\makeatother
\crefname{originalreference}{original reference}{original references}
\Crefname{originalreference}{Original reference}{Original references}
\makeatletter
\newlabel{lemma-simplicial-cocontinuous-site}{{09WD}{source}{}{sourcepointer-3}{}}
\newlabel{lemma-simplicial-cocontinuous-site@cref}{{[originalreference][0][]09WD}{[1][0][]source}{}{}{}}
\makeatother
\crefname{originalreference}{original reference}{original references}
\Crefname{originalreference}{Original reference}{Original references}
\makeatletter
\newlabel{lemma-simplicial-site-site}{{09WC}{source}{}{sourcepointer-4}{}}
\newlabel{lemma-simplicial-site-site@cref}{{[originalreference][0][]09WC}{[1][0][]source}{}{}{}}
\makeatother
\crefname{originalreference}{original reference}{original references}
\Crefname{originalreference}{Original reference}{Original references}
\makeatletter
\newlabel{remark-morphism-simplicial-sites}{{0D95}{source}{}{sourcepointer-5}{}}
\newlabel{remark-morphism-simplicial-sites@cref}{{[originalreference][0][]0D95}{[1][0][]source}{}{}{}}
\makeatother
\crefname{originallemma}{original lemma}{original lemmas}
\Crefname{originallemma}{Original lemma}{Original lemmas}
\makeatletter
\newlabel{sites-lemma-composition-cocontinuous}{{03L5}{source}{}{sourcepointer-6}{}}
\newlabel{sites-lemma-composition-cocontinuous@cref}{{[originallemma][0][]03L5}{[1][0][]source}{}{}{}}
\makeatother
\crefname{originalsection}{original section}{original sections}
\Crefname{originalsection}{Original section}{Original sections}
\makeatletter
\newlabel{sites-section-continuous-functors}{{00WU}{source}{}{sourcepointer-7}{}}
\newlabel{sites-section-continuous-functors@cref}{{[originalsection][0][]00WU}{[1][0][]source}{}{}{}}
\makeatother
\crefname{originalsection}{original section}{original sections}
\Crefname{originalsection}{Original section}{Original sections}
\makeatletter
\newlabel{sites-section-morphism-sites}{{00X0}{source}{}{sourcepointer-8}{}}
\newlabel{sites-section-morphism-sites@cref}{{[originalsection][0][]00X0}{[1][0][]source}{}{}{}}
\makeatother
\crefname{originalreference}{original reference}{original references}
\Crefname{originalreference}{Original reference}{Original references}
\makeatletter
\newlabel{situation-simplicial-site}{{09WE}{source}{}{sourcepointer-9}{}}
\newlabel{situation-simplicial-site@cref}{{[originalreference][0][]09WE}{[1][0][]source}{}{}{}}
\makeatother
\title{A morphism of simplicial sites induces a morphism of the total topoi compatible with the component inclusions}
\author{The Stacks Project Authors}
\date{stacks-project:spaces-simplicial.tex:local-snapshot-20261012; GNU FDL 1.2 or later}
\begin{document}
\maketitle
\noindent\textit{Editorial note.} All delivered mathematical text is the author's own, copied from the pinned Stacks Project chapter file sources/stacks-spaces-simplicial/pinned-original-source.tex; no statement or proof step was rewritten, reordered or removed.\par\medskip
\noindent\textit{Editorial note.} The counted proof is the author's complete argument: the two 'immediate' verifications come with their inputs in-span: (g'\_n)\textasciicircum{}\{-1\} h\_\{total,*\} = h\_\{n,*\} g\_n\textasciicircum{}\{-1\} is the section computation over objects of C'\_n described at source line 1298, and the case-B adjoint equality (lines 1340-1348) follows from the explicit shriek formulas of lemma-restriction-to-components-site (card L46149 of this batch) together with the base-change identity h\_n\textasciicircum{}\{-1\}(f'\_varphi)\textasciicircum{}\{-1\} = f\_varphi\textasciicircum{}\{-1\} h\_m\textasciicircum{}\{-1\}; the author's spellings are preserved verbatim.\par\medskip
\noindent\textit{Editorial note.} The label lemma-morphism-simplicial-sites names the counted statement, and the proof environment immediately following it in the source is delivered verbatim as its proof.\par\medskip
\noindent\textit{Editorial note.} The counted statement is carried by the author's own lemma environment in the chapter source.\par\medskip
\noindent\textit{Allowed standard prerequisites (editorial).}\begin{itemize}
\item The morphism-of-simplicial-sites datum. Same chapter, Remark \textbackslash{}label\{remark-morphism-simplicial-sites\}. Lines 1276-1278: the componentwise morphisms h\_n with their compatibility data.
\item Sites, continuous functors and the associated adjoint pair. Sites chapter, Sections \textbackslash{}ref\{sites-section-continuous-functors\} and \textbackslash{}ref\{sites-section-morphism-sites\}. Lines 1292-1296: the functors v\_\{total,s\} and v\_\{total\}\textasciicircum{}p = v\_\{total\}\textasciicircum{}s forming the adjoint pair.
\item Sheaves on the total site as systems with transition maps. Same chapter, \textbackslash{}label\{lemma-describe-sheaves-simplicial-site-site\}. Lines 1299-1304: a sheaf on C\_total is a system (F\_n, F(varphi)); used again for the exactness conclusion.
\item Sites, cocontinuous functors and composition (case B). Sites chapter, \textbackslash{}label\{sites-lemma-composition-cocontinuous\}. Lines 1334-1338: the morphism of topoi associated to the cocontinuous v\_total and the diagram commutativity.
\item Explicit description of the component shriek functors. Same chapter, \textbackslash{}label\{lemma-restriction-to-components-site\} (card L46149 of this batch). Lines 1342-1348: the case-B adjoint transfer reads off the explicit coproduct formulas.
\end{itemize}\medskip
\ifdefined\hypertarget\hypertarget{sourcepointer-1}{}\fi
\noindent\textit{Source reference.} Original reference 09WF: sources/stacks-spaces-simplicial/pinned-original-source.tex, source lines 953--953. Prerequisite pointer invoked at proof line(s) 1316; it names a standard fact of the same chapter that the counted argument uses, and every step of the counted proof built on it is retained verbatim in the delivered body.\par\medskip
\ifdefined\hypertarget\hypertarget{sourcepointer-2}{}\fi
\noindent\textit{Source reference.} Original reference 09WG: sources/stacks-spaces-simplicial/pinned-original-source.tex, source lines 970--970. Prerequisite pointer invoked at proof line(s) 1345; it names a standard fact of the same chapter that the counted argument uses, and every step of the counted proof built on it is retained verbatim in the delivered body.\par\medskip
\ifdefined\hypertarget\hypertarget{sourcepointer-3}{}\fi
\noindent\textit{Source reference.} Original reference 09WD: sources/stacks-spaces-simplicial/pinned-original-source.tex, source lines 844--844. Prerequisite pointer invoked at proof line(s) 1335; it names a standard fact of the same chapter that the counted argument uses, and every step of the counted proof built on it is retained verbatim in the delivered body.\par\medskip
\ifdefined\hypertarget\hypertarget{sourcepointer-4}{}\fi
\noindent\textit{Source reference.} Original reference 09WC: sources/stacks-spaces-simplicial/pinned-original-source.tex, source lines 797--797. Prerequisite pointer invoked at proof line(s) 1300; it names a standard fact of the same chapter that the counted argument uses, and every step of the counted proof built on it is retained verbatim in the delivered body.\par\medskip
\ifdefined\hypertarget\hypertarget{sourcepointer-5}{}\fi
\noindent\textit{Source reference.} Original reference 0D95: sources/stacks-spaces-simplicial/pinned-original-source.tex, source lines 1247--1247. Prerequisite pointer invoked at proof line(s) 1278; it names a standard fact of the same chapter that the counted argument uses, and every step of the counted proof built on it is retained verbatim in the delivered body.\par\medskip
\ifdefined\hypertarget\hypertarget{sourcepointer-6}{}\fi
\noindent\textit{Source reference.} Original lemma 03L5: sources/stacks-chow/tags.txt, source lines 4701--4701. Prerequisite pointer invoked at proof line(s) 1338; it names a standard fact of a sibling chapter that the counted argument uses, and every step of the counted proof built on it is retained verbatim in the delivered body.\par\medskip
\ifdefined\hypertarget\hypertarget{sourcepointer-7}{}\fi
\noindent\textit{Source reference.} Original section 00WU: sources/stacks-chow/tags.txt, source lines 1196--1196. Prerequisite pointer invoked at proof line(s) 1307; it names a standard fact of a sibling chapter that the counted argument uses, and every step of the counted proof built on it is retained verbatim in the delivered body.\par\medskip
\ifdefined\hypertarget\hypertarget{sourcepointer-8}{}\fi
\noindent\textit{Source reference.} Original section 00X0: sources/stacks-chow/tags.txt, source lines 1202--1202. Prerequisite pointer invoked at proof line(s) 1308; it names a standard fact of a sibling chapter that the counted argument uses, and every step of the counted proof built on it is retained verbatim in the delivered body.\par\medskip
\ifdefined\hypertarget\hypertarget{sourcepointer-9}{}\fi
\noindent\textit{Source reference.} Original reference 09WE: sources/stacks-spaces-simplicial/pinned-original-source.tex, source lines 888--888. Prerequisite pointer invoked at proof line(s) 1276; it names a standard fact of the same chapter that the counted argument uses, and every step of the counted proof built on it is retained verbatim in the delivered body.\par\medskip
\begin{lemma}
\label{lemma-morphism-simplicial-sites}
Let $\mathcal{C}_n, f_\varphi, u_\varphi$ and
$\mathcal{C}'_n, f'_\varphi, u'_\varphi$ be as in
Situation \ref{situation-simplicial-site}.
Let $h$ be a morphism between simplicial sites as in
Remark \ref{remark-morphism-simplicial-sites}.
Then we obtain a morphism of topoi
$$
h_{total} : \Sh(\mathcal{C}_{total}) \to \Sh(\mathcal{C}'_{total})
$$
and commutative diagrams
$$
\xymatrix{
\Sh(\mathcal{C}_n) \ar[d]_{g_n} \ar[r]_{h_n} &
\Sh(\mathcal{C}'_n) \ar[d]^{g'_n} \\
\Sh(\mathcal{C}_{total}) \ar[r]^{h_{total}} &
\Sh(\mathcal{C}'_{total})
}
$$
Moreover, we have $(g'_n)^{-1} \circ h_{total, *} = h_{n, *} \circ g_n^{-1}$.
\end{lemma}

\begin{proof}
Case A. Say $h_n$ corresponds to the continuous functor
$v_n : \mathcal{C}'_n \to \mathcal{C}_n$. Then we can define
a functor $v_{total} : \mathcal{C}'_{total} \to \mathcal{C}_{total}$
by using $v_n$ in degree $n$. This is clearly a continuous functor
(see definition of coverings in Lemma \ref{lemma-simplicial-site-site}).
Let
$h_{total}^{-1} = v_{total, s} :
\Sh(\mathcal{C}'_{total}) \to \Sh(\mathcal{C}_{total})$ and
$h_{total, *} = v_{total}^s = v_{total}^p :
\Sh(\mathcal{C}_{total}) \to \Sh(\mathcal{C}'_{total})$
be the adjoint pair of functors constructed and studied in
Sites, Sections \ref{sites-section-continuous-functors} and
\ref{sites-section-morphism-sites}.
To see that $h_{total}$ is a morphism of topoi
we still have to verify that $h_{total}^{-1}$ is exact.
We first observe that
$(g'_n)^{-1} \circ h_{total, *} = h_{n, *} \circ g_n^{-1}$;
this is immediate by computing sections over an object $U$
of $\mathcal{C}'_n$. Thus, if we think of a sheaf $\mathcal{F}$
on $\mathcal{C}_{total}$ as a system $(\mathcal{F}_n, \mathcal{F}(\varphi))$
as in Lemma \ref{lemma-describe-sheaves-simplicial-site-site}, then
$h_{total, *}\mathcal{F}$ corresponds to
the system $(h_{n, *}\mathcal{F}_n, h_{n, *}\mathcal{F}(\varphi))$.
Clearly, the functor
$(\mathcal{F}'_n, \mathcal{F}'(\varphi)) \to
(h_n^{-1}\mathcal{F}'_n, h_n^{-1}\mathcal{F}'(\varphi))$
is its left adjoint. By uniqueness of adjoints, we conclude that
$h_{total}^{-1}$ is given by this rule on systems. In particular,
$h_{total}^{-1}$ is exact (by the description of sheaves on
$\mathcal{C}_{total}$ given in the lemma and the exactness of
the functors $h_n^{-1}$) and we have our morphism of topoi.
Finally, we obtain $g_n^{-1} \circ h_{total}^{-1} =
h_n^{-1} \circ (g'_n)^{-1}$ as well, which proves that the
displayed diagram of the lemma commutes.

\medskip\noindent
Case B. Here we have a functor
$v_{total} : \mathcal{C}_{total} \to \mathcal{C}'_{total}$
by using $v_n$ in degree $n$. This is clearly a cocontinuous functor
(see definition of coverings in Lemma \ref{lemma-simplicial-cocontinuous-site}).
Let $h_{total}$ be the morphism of topoi associated to $v_{total}$.
The commutativity of the displayed diagram of the lemma follows
immediately from Sites, Lemma \ref{sites-lemma-composition-cocontinuous}.
Taking left adjoints the final equality of the lemma becomes
$$
h_{total}^{-1} \circ (g'_n)^{Sh}_! = g^{Sh}_{n!} \circ h_n^{-1}
$$
This follows immediately from the explicit description of the functors
$(g'_n)^{Sh}_!$ and $g^{Sh}_{n!}$ in
Lemma \ref{lemma-restriction-to-components-site},
the fact that $h_n^{-1} \circ (f'_\varphi)^{-1} =
f_\varphi^{-1} \circ h_m^{-1}$ for $\varphi : [m] \to [n]$, and
the fact that we already know $h_{total}^{-1}$ commutes
with restrictions to the degree $n$ parts of the simplicial sites.
\end{proof}

\end{document}
